\documentclass[11pt,a4paper]{article}
\usepackage{fontspec}
\usepackage{xeCJK}
\setCJKmainfont{STSong}
\setCJKsansfont{STHeiti}
\usepackage{amsmath,amssymb,amsthm}
\usepackage{mathtools}
\usepackage{bm}
\usepackage{geometry}
\usepackage{enumitem}
\usepackage{xcolor}
\usepackage[pdfencoding=auto,psdextra]{hyperref}

\geometry{margin=2.5cm}

\newcommand{\E}{\mathbb{E}}
\newcommand{\R}{\mathbb{R}}
\newcommand{\KL}{\mathrm{KL}}
\newcommand{\N}{\mathcal{N}}
\newcommand{\Xspace}{\mathcal{X}}
\newcommand{\Za}{\mathcal{Z}_{A}}
\newcommand{\Zb}{\mathcal{Z}_{B}}
\newcommand{\Enc}{\mathcal{E}}
\newcommand{\Dec}{\mathcal{D}}
\newcommand{\Tmap}{\mathsf{T}}
\newcommand{\Tinv}{\mathsf{T}^{-1}}
\newcommand{\push}{\#}
\newcommand{\Jac}{\mathbf{J}}
\newcommand{\norm}[1]{\left\lVert #1 \right\rVert}

\theoremstyle{definition}
\newtheorem{definition}{Definition}
\newtheorem{problem}{Problem}
\newtheorem{remark}{Remark}

\title{\textbf{跨隐空间速度场传输与 On-Policy 蒸馏}\\[0.3em]
\large 问题的理论形式化 \\[0.3em]
\large Cross-Latent Velocity-Field Transport for On-Policy Distillation}
\author{Flow Factory}
\date{}

\begin{document}
\maketitle

\begin{abstract}
本文只做一件事：用尽量 general 的形式化语言，把"在两个异构隐空间上的扩散
/ 流匹配模型之间做 on-policy 蒸馏"这一需求陈述为一个明确的数学问题。我们
设有两个分别在各自隐空间中去噪并生成图像的模型，目标是学习一个隐空间之间
的双向映射 $(\Tmap,\Tinv)$，使得一个模型的 velocity field 可以被
\emph{推前 / 拉回} (pushforward / pullback) 到另一个隐空间，从而让 on-policy
蒸馏的监督信号在类型与语义上都合法。本文\textbf{不涉及任何具体解法、构造或
算法比较}，只给出问题的设定、需要满足的性质，以及最终的形式化问题陈述。
\end{abstract}

\section{设定 (Setup)}
\label{sec:setup}

\paragraph{两个生成模型。}
设有一个共享的像素（图像）空间 $\Xspace\subseteq\R^{D}$，以及两个分别由各自
变分自编码器 (VAE) 诱导的隐空间
\[
\Za\subseteq\R^{d_A},\qquad \Zb\subseteq\R^{d_B},
\]
配备编码器 / 解码器对 $(\Enc_A,\Dec_A)$ 与 $(\Enc_B,\Dec_B)$，其中
$\Enc_\bullet:\Xspace\to\mathcal{Z}_\bullet$，$\Dec_\bullet:\mathcal{Z}_\bullet\to\Xspace$。
一般情形下两个隐空间是\emph{异构}的：
\[
d_A\neq d_B,\quad \Za\not\cong\Zb,\quad \Enc_A\neq\Enc_B,
\]
即通道数、空间下采样率、以及逐通道统计量都可能不同。

\paragraph{两个流 / 扩散过程。}
在每个隐空间上各有一个（流匹配 / 扩散）生成模型，由一个时间相关的
velocity field 刻画：
\[
v_A:\Za\times[0,1]\to\R^{d_A},
\qquad
v_B:\Zb\times[0,1]\to\R^{d_B}.
\]
它们各自定义了一条 probability-flow ODE，把噪声端 $t=1$ 的先验输运到数据端
$t=0$ 的隐变量边际：
\begin{equation}
\dot z_A(t) = v_A\big(z_A(t),t\big),
\qquad
\dot z_B(t) = v_B\big(z_B(t),t\big).
\label{eq:pf-ode}
\end{equation}
对应的隐变量边际记为 $p_A(\cdot,t)$ 与 $p_B(\cdot,t)$；干净端 $t=0$ 的边际
分别为 $p_A^{0}=\Enc_A\push\, p_{\mathrm{data}}$ 与 $p_B^{0}=\Enc_B\push\, p_{\mathrm{data}}$。
解码后两者生成的是同一像素分布 $p_{\mathrm{data}}$。

\paragraph{蒸馏方向。}
固定其中一个模型为 \emph{teacher}（记为 $A$，velocity $v_A$ 冻结），另一个为
\emph{student}（记为 $B$，velocity $v_B=v_\theta$，参数 $\theta$ 可训练）。蒸馏的
目标是用 teacher 在 $\Za$ 中的生成知识监督 student 在 $\Zb$ 中的生成过程。

\section{隐空间传输映射 (Latent Transport Map)}
\label{sec:transport}

由于 $v_A$ 与 $v_B$ 定义在\emph{不同}的隐空间上，teacher 的任何信号（干净
样本、velocity、转移分布）都活在 $\Za$ 中，而 student 的 rollout、loss、梯度
都活在 $\Zb$ 中。要把监督信号搬运过去，需要一个隐空间之间的双向映射。

\begin{definition}[隐空间传输映射]
\label{def:transport}
一个隐空间传输映射是一对函数
\[
\Tmap:\Za\to\Zb,\qquad \Tinv:\Zb\to\Za,
\]
满足互逆关系 $\Tinv\circ\Tmap=\mathrm{id}_{\Za}$ 与 $\Tmap\circ\Tinv=\mathrm{id}_{\Zb}$
（在异构维度 $d_A\neq d_B$ 时，互逆性只能在公共子空间 / 值域上要求，定义为
广义逆）。
\end{definition}

\begin{definition}[语义一致性 / 像素锚定]
\label{def:consistency}
称 $\Tmap$ 是\emph{像素一致}的，若它保持解码后的图像语义：对 teacher 隐变量
$z\in\Za$，
\begin{equation}
\Dec_B\big(\Tmap(z)\big)\;\approx\;\Dec_A(z),
\qquad \forall z\in\Za.
\label{eq:pixel-consistency}
\end{equation}
即 $z$ 与 $\Tmap(z)$ "代表同一张图"。这是传输具备语义正确性的判据。
\end{definition}

\begin{definition}[传输一致性误差]
\label{def:delta}
定义传输在像素空间的往返重建误差
\begin{equation}
\delta_{\Tmap}
\;=\;
\E_{z\sim p_A}\,\norm{\Dec_B\big(\Tmap(z)\big)-\Dec_A(z)},
\label{eq:delta}
\end{equation}
其中 $p_A$ 是 teacher 在相关时刻的隐变量分布。$\delta_{\Tmap}=0$ 即严格满足
像素一致 \eqref{eq:pixel-consistency}。$\delta_{\Tmap}$ 是衡量传输有损性的标量。
\end{definition}

\section{速度场的推前与拉回 (Pushforward / Pullback)}
\label{sec:pushforward}

蒸馏不仅要搬运干净样本，更要搬运\emph{整个 velocity field}，使 teacher 在
$\Za$ 中的概率流经 $\Tmap$ 之后，与 student 在 $\Zb$ 中看到的动力学一致。

\begin{problem}[velocity field 的传输]
\label{prob:pf-existence}
寻找 $(\Tmap,\Tinv)$，使得 teacher 的 probability-flow ODE
\eqref{eq:pf-ode} 在映射 $z_B(t)=\Tmap\big(z_A(t)\big)$ 下被忠实输运。
\end{problem}

形式上，若 $\Tmap$ 为 $C^1$ 微分同胚，Jacobian 记为
$\Jac_{\Tmap}(z)=\partial\Tmap/\partial z$，则由链式法则
$\dot z_B=\Jac_{\Tmap}(z_A)\,\dot z_A$，得到\emph{推前 velocity field}
\begin{equation}
(\Tmap\push v_A)(z_B,t)
\;=\;
\Jac_{\Tmap}\!\big(\Tinv(z_B)\big)\;
v_A\!\big(\Tinv(z_B),t\big),
\qquad z_B\in\Zb.
\label{eq:pushforward}
\end{equation}
对称地，把 student 空间中的状态拉回 teacher 空间需要逆映射 $\Tinv$（pullback）。
式 \eqref{eq:pushforward} 显式地表明，要把 teacher 的"意见"评估在 student
访问到的状态上，必须同时具备：
\begin{enumerate}[leftmargin=2em,itemsep=2pt]
  \item \textbf{可逆性}：能取 $\Tinv(z_B)$，以便在 student 状态 $z_B$ 处查询
        teacher；
  \item \textbf{可微性}：Jacobian $\Jac_{\Tmap}$ 存在且可计算，以便把 teacher
        velocity 正确变换到 $\Zb$。
\end{enumerate}

\section{On-Policy 蒸馏的形式化需求}
\label{sec:onpolicy}

On-policy 蒸馏的本质特征是：监督信号必须在 \emph{student 自己 rollout 出来的
状态分布} 上评估，而非在 teacher 预先采样的状态上评估。

\paragraph{student 的 on-policy 状态。}
student 用自身 velocity $v_\theta$ rollout 一条轨迹 $\{z_B(t)\}_{t}\subset\Zb$，
诱导一个被访问状态的分布 $\rho_\theta$（依赖当前参数 $\theta$，故是 on-policy 的）。

\paragraph{监督信号。}
对每个被访问状态 $z_B\sim\rho_\theta$，蒸馏要求 teacher 在该点给出"应当如何
去噪"的意见。但 teacher 的 velocity $v_A$ 只定义在 $\Za$ 上，因此其在 $\Zb$ 中
的合法取值正是\eqref{eq:pushforward}的推前场 $(\Tmap\push v_A)(z_B,t)$，它依赖
$\Tinv$（把 $z_B$ 拉回 $\Za$ 查询 teacher）与 $\Jac_{\Tmap}$（把结果变换回 $\Zb$）。

\paragraph{蒸馏目标（general 形式）。}
on-policy 蒸馏最小化 student velocity 与传输后 teacher velocity 在 student
访问分布上的偏差：
\begin{equation}
\min_{\theta}\;
\E_{t}\,\E_{z_B\sim\rho_\theta}\;
\Big[\,w(t)\,
\mathcal{M}\!\Big(v_\theta(z_B,t),\;(\Tmap\push v_A)(z_B,t)\Big)\Big],
\label{eq:distill-objective}
\end{equation}
其中 $w(t)$ 为时间权重，$\mathcal{M}(\cdot,\cdot)$ 是一个差异度量（例如二次
回归差，或由动力学诱导的转移分布散度），其具体形式不在本文范围内。该目标
良定义的前提，正是 $(\Tmap\push v_A)$ 在 $\rho_\theta$ 的支撑上可被求值——即
$\Tinv$ 与 $\Jac_{\Tmap}$ 存在且可计算。

\section{形式化问题陈述}
\label{sec:problem-statement}

综合 \S\ref{sec:setup}--\S\ref{sec:onpolicy}，跨隐空间 on-policy 蒸馏归结为以下
传输映射设计问题。

\begin{problem}[跨隐空间传输设计]
\label{prob:main}
给定两个异构隐空间 $\Za,\Zb$、对应的解码器 $\Dec_A,\Dec_B$、teacher velocity
$v_A$ 与 student velocity $v_\theta$，寻找一对映射 $(\Tmap,\Tinv)$，
$\Tmap:\Za\to\Zb$、$\Tinv:\Zb\to\Za$，使得：
\begin{enumerate}[leftmargin=2.4em,label=\textup{(P\arabic*)},itemsep=3pt]
  \item \textbf{语义一致 (semantic consistency)}：$\Tmap$ 像素一致，即传输
        一致性误差 $\delta_{\Tmap}$（Def.~\ref{def:delta}）尽可能小，
        $\Dec_B(\Tmap(z))\approx\Dec_A(z)$；
  \item \textbf{可逆 (invertibility)}：$\Tinv$ 满足 Def.~\ref{def:transport}
        的互逆关系，使 student 状态可被拉回 $\Za$ 查询 teacher；
  \item \textbf{可推前 (differentiable pushforward)}：$\Tmap$ 可微且 Jacobian
        $\Jac_{\Tmap}$ 可得，从而推前 velocity field
        \eqref{eq:pushforward} $(\Tmap\push v_A)$ 在 student 访问的状态上可
        求值。
\end{enumerate}
满足 (P1)--(P3) 时，on-policy 蒸馏目标 \eqref{eq:distill-objective} 的监督信号在
类型上合法、在语义上正确，蒸馏问题随之良定义。
\end{problem}

\begin{remark}[需求的层次性]
\label{rem:layering}
不同强度的蒸馏对 $(\Tmap,\Tinv)$ 的要求不同：仅在干净样本 $z_A^{0}$ 上做
\emph{样本传输} $z_B^{0}=\Tmap(z_A^{0})$ 只需要 (P1)；而在 student on-policy
轨迹上做 \emph{velocity field 传输} 则额外需要 (P2)--(P3)。这两个层次构成了
问题 \ref{prob:main} 的两种强度版本。
\end{remark}

\begin{remark}[范围声明]
本文只给出问题的形式化（设定、所需性质、问题陈述），不讨论 $(\Tmap,\Tinv)$
的任何具体构造、参数化、初始化或求解算法，也不对候选解法做比较与取舍。
\end{remark}

\end{document}
